\section{Impact-Driven Exploration}
\label{sec:impact_driven_exploration}

Impact-driven exploration rewards actions that produce meaningful change in the environment or in a learned representation.
The motivation is that useful exploration should not only reach unfamiliar observations, but should also involve transitions that the agent can influence.
A simple feature-space impact measure is
\begin{equation}
    I_t = \left\lVert \phi(o_{t+1})-\phi(o_t)\right\rVert_2,
\end{equation}
where $\phi$ maps observations into a representation space.
Large values indicate that the transition changed the representation substantially.

Rewarding Impact-Driven Exploration (RIDE) uses this principle to encourage agents to take actions that lead to significant feature-space changes \cite{raileanu-2020}.
RIDE combines a learned representation with an impact reward and an episodic state-count correction.
The count correction reduces the reward for repeatedly producing the same transition pattern, while the impact term encourages behavior that changes the agent's represented state.
This differs from pure novelty because the reward is tied to transition magnitude rather than only to whether a state has been visited rarely.

Impact rewards are useful in environments where the agent must learn to control movement, manipulate objects, or otherwise cause structured changes before reaching task reward.
For example, in locomotion tasks, an agent that remains stationary may avoid penalties or unstable dynamics but will not learn useful movement.
A feature-space impact term can provide pressure toward dynamic behavior even when extrinsic reward is initially weak.
In sparse-reward tasks, impact can help the agent escape local regions of the state space by valuing transitions that produce substantial state changes.

However, impact alone is not sufficient to define productive exploration.
A transition can have high impact while being harmful, random, or disconnected from task success.
Large representational changes may occur because of falls, collisions, or other unstable behavior that does not improve competence.
Furthermore, if the learned representation is poorly calibrated, the impact signal may emphasize dimensions that are easy to change but not useful for solving the task.
Impact-driven exploration therefore benefits from being combined with additional criteria that indicate whether the resulting experience supports learning.

Impact and learning progress are complementary.
Impact identifies transitions that change the representation, while learning progress identifies regions where prediction or understanding is improving.
A combined intrinsic reward can prefer transitions that are both consequential and learnable.
This distinction is important because an agent may otherwise learn to seek either low-impact predictable transitions that show small model improvement or high-impact transitions that remain unproductive.
The literature on impact-driven exploration therefore provides a key component for exploration objectives that seek to connect behavioral change with model improvement.